\documentclass[twocolumn]{article}
\usepackage{amsmath,amssymb}
\usepackage{algorithm,algorithmic}
\usepackage{booktabs,multirow}
\usepackage{subcaption,graphicx}
\usepackage{hyperref,cleveref}
\usepackage{xcolor}

\newtheorem{theorem}{Theorem}
\newcommand{\R}{\mathbb{R}}
\newcommand{\sys}{Featurd\xspace}

\title{A Feature-Exercising Note on \sys}
\author{T. Esting}

\begin{document}
\maketitle

\begin{abstract}
This document exercises algorithms, theorems, description lists, sub-figures,
citations, and cross-references. It uses \verb|inline code| too.
\end{abstract}

\section{Background}\label{sec:bg}
For $x \in \R^n$ we have $\displaystyle \sum_{i=1}^{n} x_i^2 \geq 0$, see
\eqref{eq:key}. As \citet{knuth1984} notes, typography matters
\citep[p.~5]{einstein1905}.

\begin{equation}\label{eq:key}
E = mc^2 .
\end{equation}

\begin{theorem}[Pythagoras]\label{thm:pyth}
For a right triangle, $a^2 + b^2 = c^2$.
\end{theorem}

\section{Method}\label{sec:method}
We use the following terms:
\begin{description}
\item[Latency] time to first token.
\item[\textbf{Throughput}] tokens per second.
\end{description}

\begin{algorithm}
\caption{Compute squares}\label{alg:sq}
\begin{algorithmic}[1]
\REQUIRE array $A$
\FOR{$i = 1$ to $n$}
  \STATE $A[i] \gets A[i]^2$
\ENDFOR
\RETURN $A$
\end{algorithmic}
\end{algorithm}

As \cref{alg:sq} and \Cref{thm:pyth} show, and per \cref{sec:bg}.

\begin{figure}
\begin{subfigure}{0.45\linewidth}\includegraphics{a.pdf}\caption{Left}\label{fig:l}\end{subfigure}
\begin{subfigure}{0.45\linewidth}\includegraphics{b.pdf}\caption{Right}\label{fig:r}\end{subfigure}
\caption{Two panels.}\label{fig:two}
\end{figure}

\begin{table}
\begin{tabular}{lcc}
\toprule
Method & Speed & Quality \\
\midrule
\multirow{2}{*}{Ours} & Fast & High \\
                      & Slow & Higher \\
\bottomrule
\end{tabular}
\caption{Results.}\label{tab:res}
\end{table}

Aligned derivations line up at the relation sign:
\begin{align*}
(a+b)^2 &= a^2 + 2ab + b^2 \\
        &= a^2 + b^2 + 2ab.
\end{align*}
And a single numbered aligned block:
\begin{equation}
\begin{aligned}
f(x) &= x^2 + 1 \\
g(x) &= 2x.
\end{aligned}
\end{equation}

\begin{thebibliography}{9}
\bibitem{einstein1905} A. Einstein, \emph{Zur Elektrodynamik}, 1905.
\bibitem{knuth1984} D. E. Knuth, \emph{The TeXbook}, 1984.
\end{thebibliography}

\end{document}
